% 教学假设：两对象每轮相关性均为0.5，位置注意为0.8/0.2。
\begin{tikzpicture}[
 x=1mm,y=1mm,font=\figurefont\small,
 box/.style={draw=BookBlue,line width=0.85pt,fill=BookBlue!5,
             minimum width=44mm,minimum height=13mm,align=center},
 flow/.style={->,>=latex,line width=1.1pt,draw=BookInk,
              rounded corners=1.2mm}
]
 \path[use as bounding box] (0,0) rectangle (112,65);
 \node[box] (first) at (56,56) {第一轮：A在前\\累计注意$(0.8,0.2)$};
 \node[box] (same) at (25,25) {第二轮：顺序不变\\累计注意$(1.6,0.4)$};
 \node[box,draw=BookTeal,fill=BookTeal!5] (swap) at (87,25)
    {第二轮：交换位置\\累计注意$(1,1)$};
 \draw[flow] (first.south) -- (56,41) -| (same.north);
 \draw[flow] (56,41) -| (swap.north);
 \node[text=BookMuted,font=\figurefont\footnotesize] at (56,38)
    {两轮累计相关性均为$(1,1)$};
 \node[text=BookInk] at (25,7) {累计差距$1.2$};
 \node[text=BookTeal] at (87,7) {累计差距$0$};
\end{tikzpicture}
